\subsection{Pregunta 3}

¿Por qué creen que las redes sociales nos promueven a compartir nuestra información?

\begin{respuesta}

Creemos que las empresas detras de estas redes sociales incentivan esto mismo por el simple hecho de que quieren sacarnos la mayor cantidad de información/datos posibles puesto que estos datos, desde su perspectiva, son muy valiosos y pueden ser utilizados con fines de lucro para ellos.

Con cada foto que subimos, cada ubicación que compartimos, cada amigo que etiquetamos, o cada página a la que damos "seguir", entre muchas otras interacciones que hacemos dentro de estas redes sociales, "alimentamos" un algoritmo que construye perfiles de consumidores sumamente detallados que luego pueden ser vendidos a empresas de interes (principalmente anunciantes) para que puedan mostrar publicidad dirigida que coincida exactamente con nuestros intereses y vulnerabilidades. Pues recordemos que ahora mismo vivimos en una era digital donde cada vez más personas interactúan con estos servicios, dejando un rastro de información (huella digital) que a los ojos de estas empresas es muy valiosa y que puede ser utilizada para su beneficio económico. Si bien es cierto que estas redes sociales nos ofrecen servicios gratuitos, la realidad es que nosotros somos el producto que estas empresas venden a los anunciantes, y por eso es que nos incentivan a compartir nuestra información personal.

\end{respuesta}
